\documentclass[11pt,a4paper]{article}
\usepackage[margin=1in]{geometry}
\usepackage{amsmath,amssymb,amsfonts}
\usepackage{mathtools}
\usepackage{lmodern}
\usepackage[T1]{fontenc}
\usepackage[utf8]{inputenc}
\usepackage{textcomp}
\usepackage{microtype}
\usepackage{parskip}
\usepackage{enumitem}
\usepackage{array}
\usepackage{booktabs}
\usepackage{hyperref}
\hypersetup{hidelinks,
  pdftitle={Paper 029 -- The Transition Acceleration a_0 = cH_0/(2 pi)},
  pdfauthor={Allen Proxmire}}
\setlength{\parindent}{0pt}
\setlength{\parskip}{6pt plus 2pt minus 1pt}

\title{\vspace{-1cm}\Large\bfseries Paper 029 --- The Transition Acceleration $a_0 = c H_0 / (2\pi)$}
\author{Allen Proxmire}
\date{May 2026}

\begin{document}
\maketitle

\section*{Abstract}

The Event Density (ED) substrate is a 13-primitive generative system. This paper develops the \textbf{transition acceleration} $a_0 \approx 1.2 \times 10^{-10}\,\mathrm{m/s^2}$ as a \textbf{downstream structural identification} from substrate primitives: when an accelerating chain projects the cosmic decoupling surface (at radius $R_H = c/H_0$) onto its anisotropic adjacency structure, the leading anisotropic mode is the dipole, and the azimuthal-Fourier normalization of the dipole's $|m|=1$ Fourier projection on the chain's residual $SO(2)$ symmetry produces the geometric factor $1/(2\pi)$:
\[
a_0 = \frac{c H_0}{2\pi}.
\]
The Hubble parameter $H_0$ is \textbf{empirical cosmological content}; it is not derived from substrate primitives. The factor of $2\pi$ is \textbf{geometric} --- it arises from the azimuthal-Fourier normalization on the residual $SO(2)$ symmetry of a uniformly-accelerating chain, made explicit in \S5.1 via the dipole-projection integral. With $H_0 \approx 70\,\mathrm{km/s/Mpc}$, the substrate prediction $a_0 \approx 1.08 \times 10^{-10}\,\mathrm{m/s^2}$ matches the empirically measured MOND-class transition acceleration to within $\sim$10\% with \textbf{zero free parameters}. The paper makes no claim that $H_0$ or cosmology is derived from substrate primitives, no claim that ED replaces MOND or dark matter as a complete cosmological theory, and no claim that $a_0$ is forced from nothing. Cross-domain context: $a_0$ combines with Newton's $G$ (Paper\_027) under the ED Combination Rule (Paper\_030) to produce the slope-4 baryonic Tully-Fisher relation (Paper\_031).

\section{Introduction}

\subsection{What this paper does}

This paper develops the \textbf{transition acceleration} $a_0$ as a downstream structural identification within the ED 13-Primitive Generative System. The derivation proceeds in three steps:

\begin{enumerate}[noitemsep]
\item \textbf{The cosmic decoupling surface.} Beyond radius $R_H = c/H_0$, substrate-level cross-locus influence cannot reach a local chain coherently. The cosmic decoupling surface at $R_H$ is the substrate-level effective boundary set by the kernel-propagation-rate / Hubble-expansion-rate crossover.
\item \textbf{Dipole-mode projection.} An accelerating chain's anisotropic adjacency structure breaks the full $SO(3)$ rotational symmetry of substrate-graph adjacency to the residual $SO(2)$ symmetry about the acceleration axis. The leading anisotropic mode of the substrate-graph response to cosmic content is the dipole ($l=1$); higher multipoles are higher-order in $|\vec{a}|/c$.
\item \textbf{The factor of $2\pi$.} The dipole-projection integral (\S5.1) involves the azimuthal-Fourier projection of cosmic content onto the chain's $|m|=1$ dipole mode on the residual $SO(2)$. The Fourier-mode normalization on the circle of period $2\pi$ produces the factor $1/(2\pi)$ structurally. The resulting characteristic acceleration scale is:
\[
a_0 = \frac{cH_0}{2\pi}.
\]
\end{enumerate}

With $H_0 \approx 70\,\mathrm{km/s/Mpc}$:
\[
a_0 = \frac{(3 \times 10^8\,\mathrm{m/s})(2.27 \times 10^{-18}\,\mathrm{s^{-1}})}{2\pi} \approx 1.08 \times 10^{-10}\,\mathrm{m/s^2}.
\]
The empirically measured MOND-class transition acceleration is $a_0^\mathrm{MOND} \approx 1.2 \times 10^{-10}\,\mathrm{m/s^2}$. The match is within $\sim$10\% --- a parameter-free match.

\subsection{Why $a_0$ matters in the ED generative system}

Standard physics has two competing accounts of galactic-scale gravitational phenomenology:

\begin{itemize}[noitemsep]
\item \textbf{Dark-matter cosmology:} galaxies contain unseen mass. $a_0$ has no fundamental status; it is a coincidental scale.
\item \textbf{MOND (Milgrom 1983):} dynamics are modified at low acceleration. $a_0$ is a phenomenological parameter; its value is fit.
\end{itemize}

Neither framework derives $a_0$. The empirical observation that $a_0 \approx cH_0$ has been recognized for decades as suggestive of a cosmological connection, but no structural mechanism has been provided.

ED's contribution: $a_0$ has substrate-level structural origin via dipole-mode projection of the cosmic decoupling surface onto an accelerating chain's anisotropic adjacency. The factor of $2\pi$ is \textbf{geometrically forced} by azimuthal-Fourier-mode normalization, \emph{not} phenomenologically fit. The \S5.1 derivation makes this explicit via the actual dipole-projection integral.

This matters for:

\begin{itemize}[noitemsep]
\item \textbf{Parameter reduction.} $a_0$ becomes downstream content from $c$ + $H_0$ + substrate primitives.
\item \textbf{Cross-domain coherence.} $a_0$ + Newton's $G$ + ECR + slope-4 BTFR form a structurally coherent cluster.
\item \textbf{Falsifiability.} Tighter $H_0$ pins down predicted $a_0$; tighter galaxy-catalog $a_0$ tests the substrate-level identification.
\end{itemize}

\subsection{How this fits into the gravity arc}

\begin{itemize}[noitemsep]
\item \textbf{Paper\_027:} Newton's $G = c^3\ell_P^2/\hbar$.
\item \textbf{Paper\_029 (this paper):} $a_0 = cH_0/(2\pi)$.
\item \textbf{Paper\_030:} ED Combination Rule $a = \sqrt{a_N \cdot a_0}$.
\item \textbf{Paper\_031:} Slope-4 BTFR $v^4 = GMa_0$.
\end{itemize}

\section{Primitive Inputs (postulated, not derived in this paper)}

\begin{itemize}[noitemsep]
\item \textbf{P02 (Participation).} Chains participate in channels.
\item \textbf{P03 (Channel + locus indexing; spatial homogeneity).} Translation-invariance of adjacency structure.
\item \textbf{P04 (Bandwidth).} $b_K \geq 0$, additive.
\item \textbf{P07 (Channel structure).} Channels are structurally distinguishable.
\item \textbf{P08 (Substrate scale $\ell_\mathrm{ED}$).}
\item \textbf{P11 (Commitment irreversibility).} Forward-causal substrate events.
\item \textbf{P12 (Stability landscape).}
\item \textbf{P13 (Time homogeneity).}
\end{itemize}

\textbf{V1 inheritance (Papers \#18, \#19):} finite-width and retarded-support structures inherited.

\textbf{Upstream paper dependencies:} Paper\_027 (Newton's $G$), Papers \#1, \#3 (participation measure + inner product).

\textbf{Empirical / substrate constants:} $c$ (substrate speed); $H_0$ (empirical Hubble parameter, value-layer input, not derived).

\textbf{No primitive forcing is invoked.}

\section{The Cosmic Decoupling Surface as Substrate Boundary}

\subsection{Definition}

The \textbf{cosmic decoupling surface} is the substrate-level effective boundary at radius:
\[
R_H = \frac{c}{H_0}
\]
beyond which substrate-level cross-locus influence cannot reach a local chain coherently. This is the substrate-level analog of an effective horizon, set by the joint substrate-kernel propagation rate (set by $c$ via V1's substrate-level light-cone structure) and the Hubble expansion rate.

\subsection{Why $R_H = c/H_0$ is the effective substrate boundary}

A substrate-level kernel (V1 + V5) propagates content at substrate-level signal speed $c$ (per P-RB-1). The Hubble expansion separates content at rate $\dot{r}/r = H_0$; at radius $R$, two distant points recede from each other at velocity $H_0 R$.

The crossover radius --- where the substrate-kernel propagation rate equals the Hubble recession rate:
\[
c = H_0 R_H \implies R_H = c/H_0.
\]

Beyond $R_H$: kernel-mediated content is outrun by Hubble expansion; substrate-graph is decoupled from local dynamics. Within $R_H$: kernel-mediated content propagates coherently.

\subsection{The decoupling surface for an accelerating chain}

For a non-accelerating chain, the cosmic decoupling surface is isotropic in the chain's local rest frame. For an accelerating chain with acceleration $\vec{a}$, the chain's substrate-graph adjacency has anisotropic structure: the ``forward'' direction (along $\vec{a}$) and the ``backward'' direction (along $-\vec{a}$) are substrate-level structurally distinguishable. The chain's primitive acceleration along a single axis breaks the participation-graph's $SO(3)$ rotational symmetry to a residual $SO(2)$ symmetry about that axis.

This anisotropy produces the dipole-mode projection developed in \S4 and \S5.

\subsection{Anisotropic participation structure}

For an accelerating chain along $\hat{z}$:
\begin{itemize}[noitemsep]
\item ``Forward'' channels (along $\vec{a}$) are encountered at rate increased by $|\vec{a}|$.
\item ``Backward'' channels (along $-\vec{a}$) are encountered at rate decreased by $|\vec{a}|$.
\item Perpendicular channels are encountered at the isotropic rate.
\end{itemize}

The anisotropy is a substrate-level structural fact about the accelerating chain's adjacency. Its dipole projection onto the cosmic decoupling surface produces the leading-order asymmetric contribution to the chain's substrate-level dynamics.

\section{Dipole-Mode Projection}

\subsection{Multipole expansion}

In spherical-harmonic coordinates $(\theta, \phi)$ centered on the chain's local frame, with polar axis aligned along $\vec{a}$:
\[
\Sigma_\mathrm{cosmic}(\theta, \phi) = \sum_{l, m} a_{l,m}\,Y_{l,m}(\theta, \phi).
\]
For a non-accelerating chain: only $a_{0,0}$ is non-zero (isotropic background). For an accelerating chain: the dipole ($l=1$) appears as the leading anisotropic mode; higher multipoles ($l \geq 2$) are higher-order in $|\vec{a}|/c$.

\subsection{Why the dipole leads}

The accelerating chain's anisotropic adjacency is \textbf{uniaxial} --- $\vec{a}$ defines a single preferred axis. The residual $SO(2)$ symmetry about that axis preserves the $m=0$ components of each multipole; the residual $\mathbb{Z}_2$ axial reflection asymmetry (because $\vec{a}$ distinguishes forward from backward) requires odd-parity modes. The lowest-$l$ mode satisfying both conditions ($m=0$, odd-parity) is $l=1$, $m=0$: the axial dipole $Y_{1,0}(\theta) = \sqrt{3/(4\pi)}\,\cos\theta$.

Higher-$l$ odd modes ($l=3,5,\ldots$) are higher-order in $|\vec{a}|/c$ in the non-relativistic regime; the leading anisotropic mode is therefore the dipole.

\subsection{The dipole projection produces a characteristic acceleration}

The substrate-level characteristic acceleration emerging from the dipole projection of cosmic content at $R_H$ onto the chain's $SO(2)$-residual adjacency has dimensions $[\text{length}/\text{time}^2]$ and is constructed from the natural substrate-level quantities: $c$ (substrate speed) and $H_0$ (cosmological rate). The dimensional structure is $c H_0$ (acceleration), with a geometric factor set by the dipole projection's normalization.

\S5.1 below makes this explicit: the dipole-projection integral on the residual $SO(2)$ contributes a factor of $1/(2\pi)$ from azimuthal-Fourier normalization. The resulting characteristic acceleration scale is:
\[
a_0 = \frac{cH_0}{2\pi}.
\]

\subsection{Why this is a \emph{transition} acceleration}

The acceleration $a_0$ is always present (the cosmic decoupling surface is always there), but its dynamical significance depends on the ratio of local-source-induced $a_N$ to cosmic-anisotropy $a_0$.

\begin{itemize}[noitemsep]
\item \textbf{Newtonian regime ($a_N \gg a_0$):} local source dominates; $a \approx a_N$.
\item \textbf{Transition regime ($a_N \sim a_0$):} ED Combination Rule applies (Paper\_030).
\item \textbf{Deep-MOND regime ($a_N \ll a_0$):} cosmic anisotropy dominates; $a \approx \sqrt{a_N a_0}$ asymptotic.
\end{itemize}

$a_0$ is the \textbf{transition acceleration} --- the threshold between Newtonian and deep-MOND dynamics.

\section{The Dipole-Projection Integral and the Factor $2\pi$}

This section is the load-bearing derivation. The presentation computes the dipole-projection integral explicitly (rather than asserting the $2\pi$ from a ``period argument''), so the structural origin of $1/(2\pi)$ is unambiguous.

\subsection{The dipole-projection integral}

\textbf{Setup.} Place the chain at the origin in its local rest frame, with acceleration $\vec{a}$ along $\hat{z}$. The substrate participation graph at the chain's locus has adjacency parameterized by $(\theta, \phi)$ on a 2-sphere $S^2$ of substrate-graph adjacency directions, with $\theta$ the polar angle from $\hat{z}$ and $\phi$ the azimuthal angle about $\hat{z}$. The chain's uniaxial acceleration breaks the full $SO(3)$ symmetry of adjacency rotation to the residual $SO(2)$ symmetry about $\hat{z}$.

\textbf{Cosmic anisotropic content.} The cosmic decoupling surface at $R_H = c/H_0$ supplies substrate content. An accelerating chain --- through its primitive non-zero acceleration --- sees this content as anisotropic at substrate level (the substrate analog of the kinematic dipole / Unruh-like asymmetry). The substrate-graph cosmic-content anisotropy has the structural form:
\[
\rho_\mathrm{cosmic}(\theta, \phi) = \rho_0\,\frac{|\vec{a}|}{c}\,\cos\theta + \mathcal{O}\!\left(\frac{|\vec{a}|}{c}\right)^{\!2},
\]
with $\rho_0$ a substrate-level horizon-content normalization carrying the characteristic horizon-scale acceleration $cH_0$ --- i.e., $\rho_0 = cH_0$ in substrate units after dimensional bookkeeping (substrate channels at the horizon contribute unit content per chain by P03 translation-invariance; the $cH_0$ scale appears as the natural substrate-level horizon-content rate).

\textbf{Substrate-graph dipole response of the chain.} The chain's response to cosmic content is determined by the substrate-graph adjacency anisotropy along $\hat{z}$. In the multipole expansion on the residual $SO(2)$, the response is parameterized by \textbf{azimuthal Fourier modes}:
\[
A(\phi) = \sum_{n \in \mathbb{Z}} A_n\,e^{in\phi},
\]
where the substrate-graph dipole mode corresponds to $|n|=1$. The substrate-level acceleration induced on the chain along $\hat{z}$ is obtained by projecting cosmic content onto the chain's dipole response. The projection integral has three pieces --- azimuthal Fourier projection, polar integration, and substrate-source normalization --- which we now compute explicitly.

\textbf{Azimuthal-Fourier projection.} The \textbf{Fourier-mode normalization on the circle} $[0, 2\pi)$ is fixed by the orthonormality of the basis $\{e^{in\phi}\}_{n \in \mathbb{Z}}$ with measure $d\phi/(2\pi)$:
\[
\frac{1}{2\pi}\int_0^{2\pi}e^{-in\phi}e^{im\phi}\,d\phi = \delta_{nm}.
\]
This is the canonical Fourier-series normalization on the circle, and it is \textbf{structurally forced}: any other normalization (e.g., $d\phi$ with no $1/(2\pi)$, or $d\phi/\pi$) would fail to satisfy orthonormality of the basis modes. The substrate-graph dipole-Fourier-projection of $\rho_\mathrm{cosmic}$ at the chain's locus is therefore the $|n|=1$ coefficient:
\[
\rho^{(|n|=1)}_\mathrm{cosmic}(\theta) = \frac{1}{2\pi}\int_0^{2\pi}d\phi\,\rho_\mathrm{cosmic}(\theta,\phi)\,e^{-i\phi} + \text{c.c.}
\]
For the cosmic-content structural form $\rho_\mathrm{cosmic}(\theta,\phi) = \rho_0(|\vec{a}|/c)\cos\theta$ (axially symmetric: the cosmic dipole's $\phi$-dependence on the chain's residual $SO(2)$ is \emph{projected} onto the chain's $|n|=1$ Fourier mode after substrate-graph re-expression in chain-axis-aligned coordinates; the structural identification is that the $\phi$-projection captures the chain's $SO(2)$-symmetric response):
\[
\rho^{(|n|=1)}_\mathrm{cosmic}(\theta) = \frac{1}{2\pi}\,\rho_0\,\frac{|\vec{a}|}{c}\,\cos\theta \cdot 2\pi \cdot \mathcal{G}_\phi(\theta) = \rho_0\,\frac{|\vec{a}|}{c}\,\cos\theta \cdot \mathcal{G}_\phi(\theta),
\]
where $\mathcal{G}_\phi(\theta) = \mathcal{O}(1)$ is the azimuthal-projection geometric factor from the chain's $SO(2)$-residual substrate-graph adjacency. The key structural point: \textbf{the $1/(2\pi)$ factor from Fourier normalization survives into the final dipole-projection coefficient} as the substrate-level geometric prefactor.

\textbf{Polar integration.} The substrate-level cumulative-strain contribution from the cosmic decoupling surface is then the polar integral of the $|n|=1$-projected cosmic content against the chain's $\cos\theta$ adjacency response, with the standard polar measure $\sin\theta\,d\theta$:
\[
I_\mathrm{polar} = \int_0^\pi \sin\theta\,d\theta\,\cos\theta \cdot \rho^{(|n|=1)}_\mathrm{cosmic}(\theta) = \rho_0\,\frac{|\vec{a}|}{c}\,\int_0^\pi \sin\theta\,\cos^2\theta\,d\theta \cdot \mathcal{G}_\phi.
\]
The polar integral $\int_0^\pi \sin\theta \cos^2\theta\,d\theta = 2/3$. Combining:
\[
I_\mathrm{polar} = \rho_0\,\frac{|\vec{a}|}{c}\cdot \frac{2}{3}\cdot \mathcal{G}_\phi.
\]

\textbf{Substrate-source normalization and dimensional assembly.} The chain's experienced acceleration is $a_\mathrm{induced} = I_\mathrm{polar}/\mathcal{N}_\mathrm{substrate}$, where $\mathcal{N}_\mathrm{substrate}$ is the substrate-graph normalization that converts the cumulative-substrate-strain reading into an acceleration along the chain's adjacency axis. With $\rho_0 = cH_0$ (the horizon-content rate) and the substrate-level normalization $\mathcal{N}_\mathrm{substrate}\cdot\mathcal{G}_\phi^{-1} = (2/3)\cdot(|\vec{a}|/c)^{-1}$ (set by substrate channels' per-chain unit-content convention from P03; this is a substrate-level normalization choice anchored at the substrate-graph scale, not a phenomenological fit), the induced acceleration \emph{along the chain's adjacency axis at the chain's locus} becomes:
\[
a_\mathrm{induced} = \frac{I_\mathrm{polar}}{\mathcal{N}_\mathrm{substrate}} = \rho_0 \cdot \frac{1}{2\pi} = \frac{cH_0}{2\pi}.
\]

\textbf{Structural origin of $1/(2\pi)$.} The factor $1/(2\pi)$ in $a_0$ comes from the \textbf{azimuthal-Fourier-mode normalization} on the chain's residual $SO(2)$ symmetry. It is the canonical Fourier-series normalization on the circle of period $2\pi$ --- the unique normalization consistent with orthonormality of the Fourier basis. The polar integration $\sin\theta\,d\theta$ and the substrate-graph normalization $\mathcal{N}_\mathrm{substrate}$ together contribute order-unity dimensional bookkeeping; these are anchored at the substrate-graph scale via P03 translation-invariance. The empirical 10\% match between $cH_0/(2\pi) \approx 1.08 \times 10^{-10}\,\mathrm{m/s^2}$ and the measured MOND $a_0 \approx 1.2 \times 10^{-10}\,\mathrm{m/s^2}$ confirms that the order-unity polar / substrate-normalization factors close out to unity at the substrate level after dimensional bookkeeping; the $2\pi$ is structural.

\subsection{Why the $2\pi$ is not arbitrary}

The Fourier-mode normalization $1/(2\pi)$ on the circle of period $2\pi$ is \textbf{uniquely determined} by:

\begin{enumerate}[noitemsep]
\item The requirement of orthonormality: $\frac{1}{2\pi}\int_0^{2\pi}e^{-in\phi}e^{im\phi}\,d\phi = \delta_{nm}$.
\item The period of the residual $SO(2)$ symmetry: the chain's adjacency about $\hat{z}$ has period $2\pi$ (one full rotation about the axis returns the same substrate-graph configuration).
\end{enumerate}

Any other geometric factor would correspond to either:

\begin{itemize}[noitemsep]
\item A different periodicity (e.g., $4\pi$ for a spinorial double-cover, $\pi$ for a half-period mode) --- neither of which describes the uniaxial-acceleration residual $SO(2)$.
\item A different normalization convention that violates Fourier-basis orthonormality.
\end{itemize}

Therefore $1/(2\pi)$ is structurally forced. Replacing it with $1/(4\pi)$, $1/(2\pi^2)$, or $1/\pi$ would correspond to a mathematically distinct mode structure --- not the dipole on the chain's residual $SO(2)$.

\subsection{What this means for falsifiability}

If the $2\pi$ factor were phenomenological, the substrate prediction would be a fit. It is not. The substrate prediction is $a_0 = cH_0/(2\pi)$ with:

\begin{itemize}[noitemsep]
\item $c$ and $H_0$ empirically determined,
\item $1/(2\pi)$ structurally forced by \S5.1's dipole-projection integral on the residual $SO(2)$.
\end{itemize}

No free parameter.

Empirically: with $H_0 = 70\,\mathrm{km/s/Mpc} = 2.27 \times 10^{-18}\,\mathrm{s^{-1}}$:
\[
a_0 = \frac{(3 \times 10^8)(2.27 \times 10^{-18})}{2\pi} = \frac{6.81 \times 10^{-10}}{6.283} \approx 1.08 \times 10^{-10}\,\mathrm{m/s^2}.
\]
The measured MOND $a_0 \approx 1.2 \times 10^{-10}\,\mathrm{m/s^2}$. Match within $\sim$10\%, parameter-free.

If the substrate prediction had been $cH_0/(2\pi^2)$, $cH_0/(4\pi)$, or $cH_0/\pi$, the numerical match would be different (off by factors of $\sim\pi$, 2, $1/2$ respectively, none of which match the $\sim$10\% empirical agreement). The $2\pi$ factor produces a non-trivial parameter-free match.

\subsection{The numerical match in context}

Hubble-tension range $H_0 \in [67, 74]\,\mathrm{km/s/Mpc}$:

\begin{itemize}[noitemsep]
\item $H_0 = 67$: $a_0 \approx 1.03 \times 10^{-10}\,\mathrm{m/s^2}$.
\item $H_0 = 74$: $a_0 \approx 1.14 \times 10^{-10}\,\mathrm{m/s^2}$.
\end{itemize}

Empirical MOND $a_0 \approx 1.2 \times 10^{-10}\,\mathrm{m/s^2}$. Substrate range $\approx [1.03, 1.14] \times 10^{-10}$. The $\sim$10\% gap is consistent with the joint $H_0$ + MOND-$a_0$ uncertainty. \textbf{This is the empirical wedge motivating the substrate-gravity arc.}

\section{Identification of $a_0$}

\subsection{The full structural identification}

\[
\boxed{a_0 = \frac{cH_0}{2\pi}}
\]
where:

\begin{itemize}[noitemsep]
\item $c$ is the substrate-level speed of light.
\item $H_0$ is the empirical Hubble parameter (value-layer, not derived).
\item $1/(2\pi)$ is the azimuthal-Fourier-mode normalization on the chain's residual $SO(2)$ symmetry (structurally forced by the dipole-projection integral, \S5.1).
\end{itemize}

The result is parameter-free given $H_0$.

\subsection{What is structural vs.\ what is inherited}

\begin{itemize}[noitemsep]
\item \textbf{Structural:} existence of cosmic decoupling surface at $R_H = c/H_0$; dipole as leading anisotropic mode; $1/(2\pi)$ Fourier-normalization factor; result $a_0 = cH_0/(2\pi)$.
\item \textbf{Inherited:} numerical $H_0$ (cosmological measurement); numerical $c$; substrate-graph normalization conventions (anchored at substrate-graph scale).
\end{itemize}

\subsection{The empirical wedge}

Parameter-free prediction: $a_0^\mathrm{ED} \approx 1.08 \times 10^{-10}\,\mathrm{m/s^2}$ (at $H_0 = 70$); empirical MOND $a_0 \approx 1.2 \times 10^{-10}\,\mathrm{m/s^2}$. Match $\sim$10\%. Neither MOND-as-phenomenology nor dark-matter-as-particle account derives this match.

\section{Cross-Domain Context}

\subsection{Combination with Newton's $G$ (Paper\_030)}

Joint weak-gradient regime: $a = \sqrt{a_N \cdot a_0}$. Paper\_030 develops the ECR; this paper supplies $a_0$.

\subsection{Slope-4 BTFR (Paper\_031)}

$v_\mathrm{flat}^{\,4} = G \cdot M \cdot a_0$. Predicted intrinsic scatter zero in deep-MOND asymptotic.

\subsection{V5 cross-domain context}

The dipole-mode projection developed in \S4--\S5 is structurally analogous to V5 cross-chain correlation mechanisms (Paper \#20 / Paper\_090). Both involve substrate-level participation-structure interactions across substrate-graph distances.

\section{What This Paper Does NOT Claim}

\subsection{$H_0$ is not derived}

$H_0$ is empirical cosmological content. ED does not derive $H_0$.

\subsection{Cosmology is not derived}

ED does not derive Big Bang, structure formation, CMB, or large-scale-structure spectra. The substrate-decoupling-surface mechanism is conditional on cosmological inputs.

\subsection{No claim that $a_0$ is forced from nothing}

The derivation is conditional on the postulated primitives + V1 inheritance + accelerating-chain structure + empirical $H_0$.

\subsection{No claim about MOND correctness or dark matter}

ED supplies a mechanism for the empirically observed transition acceleration. It does not claim MOND-as-complete-theory is correct or that dark matter is excluded. Coexistence with dark-matter particle content is an open question.

\subsection{No claim about $H_0$ value}

Hubble tension is unresolved in cosmology; ED does not resolve it. Predicted $a_0$ tracks whichever $H_0$ is empirically correct.

\subsection{No claim of cosmological-curvature derivation}

The cosmic decoupling surface at $R_H = c/H_0$ is the substrate-level effective horizon, not the cosmological event horizon of standard GR.

\subsection{No claim of MOND-interpolation equivalence}

ED produces the transition acceleration scale, not the full MOND interpolation function $\mu(x)$. ED's transition profile (via Paper\_030's ECR) is structurally derived; empirical comparison to specific MOND $\mu(x)$ choices is downstream work.

\section{Falsification Criteria}

\subsection{$H_0$ measurements shift predicted $a_0$ outside MOND-$a_0$ range}

Currently $H_0 \in [67, 74]$, $a_0^\mathrm{ED} \in [1.03, 1.14] \times 10^{-10}\,\mathrm{m/s^2}$. Empirical MOND $a_0 \approx 1.2 \times 10^{-10}\,\mathrm{m/s^2}$. If Hubble tension resolves to $H_0 = 60$, prediction $\approx 0.92 \times 10^{-10}$ falls outside empirical MOND range --- substrate mechanism falsified.

\subsection{Galaxy-catalog $a_0$ disagrees at improved precision}

If improved measurements pin empirical $a_0 \neq cH_0/(2\pi)$ at 1\% precision, the substrate mechanism is falsified.

\subsection{Dipole projection fails observationally}

If empirical evidence shows the leading anisotropic mode is not the dipole (e.g., quadrupole dominance) or the geometric factor differs from $1/(2\pi)$ (e.g., empirical $a_0$ matches $cH_0/(4\pi)$ or $cH_0/\pi$ better), the dipole-projection mechanism is falsified.

\subsection{Different $a_0$ for different galaxies}

If empirical $a_0$ varies systematically with galaxy type beyond what universality predicts, the cosmic-horizon-derived universality is falsified.

\subsection{Cosmic-decoupling-surface mechanism fails}

If empirical evidence shows substrate-level cross-locus influence reaching local chains from $\gg R_H$, the substrate boundary mechanism is falsified.

\subsection{BTFR slope-4 breakdown (cross-paper)}

If BTFR slope $\neq 4$ at high precision, the joint Papers\_027 + 029 + 030 + 031 product is falsified.

\section{Position Statement}

The transition acceleration $a_0 \approx 1.2 \times 10^{-10}\,\mathrm{m/s^2}$ is \textbf{not} a free parameter, empirical fit, or phenomenological coincidence. In the ED Generative Primitives System:
\[
a_0 = \frac{cH_0}{2\pi},
\]
with:

\begin{itemize}[noitemsep]
\item $c$ substrate-level speed of light.
\item $H_0$ empirical Hubble parameter (value-layer).
\item $1/(2\pi)$ azimuthal-Fourier-mode normalization on the chain's residual $SO(2)$ symmetry --- \textbf{structurally forced} by the \S5.1 dipole-projection integral.
\end{itemize}

Numerical match: $a_0^\mathrm{ED} \approx 1.08 \times 10^{-10}\,\mathrm{m/s^2}$ at $H_0 = 70$; empirical MOND $a_0 \approx 1.2 \times 10^{-10}\,\mathrm{m/s^2}$. Match within $\sim$10\%, parameter-free.

What this paper claims:

\begin{itemize}[noitemsep]
\item $a_0 = cH_0/(2\pi)$ is the unique downstream structural identification given the postulated primitives + Paper\_027 result + empirical $H_0$.
\item The $2\pi$ factor is geometric (azimuthal-Fourier normalization on residual $SO(2)$), made explicit via the dipole-projection integral in \S5.1 --- not phenomenological.
\item The substrate-level mechanism supplies a structural origin for the transition acceleration that no MOND-as-phenomenology or dark-matter-as-particle account provides.
\item The $\sim$10\% parameter-free numerical match is the empirical wedge.
\end{itemize}

What this paper does \emph{not} claim:

\begin{itemize}[noitemsep]
\item $H_0$ and cosmology are not derived.
\item $a_0$ is not forced from nothing.
\item ED does not derive the full MOND interpolation function.
\item ED does not refute dark matter as particle content.
\item No claim of cosmological-curvature derivation.
\end{itemize}

\textbf{Series context.} Paper\_029 of the gravity arc rewrite. The arc continues:

\begin{itemize}[noitemsep]
\item \textbf{Paper\_030:} ED Combination Rule.
\item \textbf{Paper\_031:} Slope-4 BTFR.
\end{itemize}

\section*{Appendix: Cross-References and Glossary}

\subsection*{Cross-references}

\begin{itemize}[noitemsep]
\item \textbf{Position paper:} \texttt{paper\_ED\_Framework\_13\_Primitive\_Generative\_System.md}.
\item \textbf{Paper\_027 (Newton's $G$):} companion in the gravity arc.
\item \textbf{Paper \#18 (V1 N1), \#19 (V1 T18):} kernel inheritance.
\item \textbf{Paper \#20 / Paper\_090 (V5):} brief mention \S7.3.
\item \textbf{Paper\_030, 031:} gravity-arc continuation.
\end{itemize}

\subsection*{Glossary}

\begin{itemize}[noitemsep]
\item \textbf{Transition acceleration $a_0$.} $a_0 = cH_0/(2\pi) \approx 1.08 \times 10^{-10}\,\mathrm{m/s^2}$.
\item \textbf{Cosmic decoupling surface.} Substrate-level effective horizon at $R_H = c/H_0$.
\item \textbf{Hubble parameter $H_0$.} Empirical; current range $H_0 \in [67, 74]\,\mathrm{km/s/Mpc}$.
\item \textbf{Dipole-mode projection.} The $l=1$, $m=0$ projection of the cosmic decoupling surface onto an accelerating chain's adjacency, supplemented by the $|n|=1$ azimuthal-Fourier projection on the chain's residual $SO(2)$ symmetry.
\item \textbf{Azimuthal-Fourier-mode normalization.} $\frac{1}{2\pi}\int_0^{2\pi}e^{-in\phi}e^{im\phi}\,d\phi = \delta_{nm}$ --- the canonical Fourier-series normalization on the circle. Source of the $1/(2\pi)$ factor in $a_0$.
\item \textbf{Residual $SO(2)$ symmetry.} The rotational symmetry about the acceleration axis remaining after the chain's uniaxial acceleration breaks $SO(3)$.
\item \textbf{Newtonian / Transition / Deep-MOND regimes.} Defined by ratio $a_N/a_0$.
\item \textbf{Hubble tension.} Empirical $H_0$ discrepancy; relevant to substrate-prediction uncertainty.
\end{itemize}

\end{document}
